\documentclass[11pt, a4paper]{article}

%------------------------------------------------------------------------------
% PACKAGES
%------------------------------------------------------------------------------
\usepackage[a4paper, top=1.2cm, bottom=1.2cm, left=1.4cm, right=1.4cm]{geometry}
\usepackage[T1]{fontenc}
\usepackage[utf8]{inputenc}
\usepackage{lmodern}
\usepackage{microtype}
\usepackage{parskip}
\usepackage{enumitem}
\usepackage{titlesec}
\usepackage{xcolor}
\usepackage{tabularx}
\usepackage{array}
\usepackage{booktabs}
\usepackage{hyperref}
\usepackage{fontawesome5}

%------------------------------------------------------------------------------
% FARBEN
%------------------------------------------------------------------------------
\definecolor{headingcolor}{RGB}{0,0,0}
\definecolor{linkcolor}{RGB}{26, 90, 150}
\definecolor{rulecolor}{RGB}{180,180,180}

%------------------------------------------------------------------------------
% HYPERREF
%------------------------------------------------------------------------------
\hypersetup{
    colorlinks = true,
    urlcolor   = linkcolor,
    linkcolor  = linkcolor,
    pdftitle   = {Lebenslauf -- Pravin Mani Kannan},
    pdfauthor  = {Pravin Mani Kannan},
}

%------------------------------------------------------------------------------
% ABSCHNITTSFORMATIERUNG
%------------------------------------------------------------------------------
\titleformat{\section}
  {\large\bfseries\color{headingcolor}}
  {}
  {0em}
  {}
  [\vspace{-17pt}\textcolor{rulecolor}{\rule{\linewidth}{0.5pt}}]

\titlespacing{\section}{0pt}{4pt}{0pt}

%------------------------------------------------------------------------------
% EIGENE BEFEHLE
%------------------------------------------------------------------------------

% Eintrag mit Titel links und Datum/Ort rechts
\newcommand{\cventry}[3]{%
  \noindent
  \begin{tabularx}{\linewidth}{@{}X r@{}}
    \textbf{#1} & \small #3
  \end{tabularx}%
  \ifx&#2&\else\vspace{-\parskip}\vspace{0pt}\par\noindent\textit{#2}\fi
  \vspace{1pt}
}

% Projekteintrag (fett + kursiver Tag, keine rechte Spalte)
\newcommand{\project}[2]{%
  \noindent\textbf{#1} \textit{[#2]}%
  \vspace{-1pt}
}

% Projekteintrag mit Datum rechts
\newcommand{\projectdate}[3]{%
  \noindent
  \begin{tabularx}{\linewidth}{@{}X r@{}}
    \textbf{#1} \textit{[#2]} & \small #3
  \end{tabularx}%
  \vspace{1pt}
}

% Aufzählungsliste mit engem Abstand
\newenvironment{cvlist}{%
  \begin{itemize}[leftmargin=1.2em, itemsep=-1pt, parsep=1pt, topsep=-1pt]
}{%
  \end{itemize}
}

% Kenntniszeile: fettes Label + normaler Text
\newcommand{\skillrow}[2]{%
  \noindent\textbf{#1:} #2\par\vspace{0pt}
}

%----------------------------------------------------------
% PERSÖNLICHE DATEN: zweispaltige Label/Wert-Tabelle
%
% \infolabelwidth  -- Breite der Label-Spalte (justiertbar)
% \inforowgap      -- Abstand zwischen den Zeilen (justiertbar)
%----------------------------------------------------------
\newlength{\infolabelwidth}
\setlength{\infolabelwidth}{4.5cm}   % <-- Labelbreite anpassen

\newlength{\inforowgap}
\setlength{\inforowgap}{-2pt}         % <-- Zeilenabstand anpassen

\newcommand{\inforow}[2]{%
  \noindent
  \begin{tabularx}{\linewidth}{@{}p{\infolabelwidth} X@{}}
    \textbf{#1} & #2
  \end{tabularx}\par\vspace{\inforowgap}
}

%------------------------------------------------------------------------------
% DOKUMENT
%------------------------------------------------------------------------------
\begin{document}
\pagestyle{empty}

%==============================================================================
% KOPFZEILE
%==============================================================================
\begin{center}
  {\Huge\bfseries PRAVIN MANI KANNAN}\\[3pt]
  \faEnvelope\ \href{mailto:pravin.mani@fau.de}{E-Mail}
  \quad\textbar\quad
  \faLinkedin\ \href{http://www.linkedin.com/in/mpravin21/}{LinkedIn}
  \quad\textbar\quad
  \faGithub\ \href{https://github.com/Usernotfoundoops}{GitHub}
\end{center}

\vspace{4pt}

%==============================================================================
% PERSÖNLICHE DATEN
%==============================================================================
\section{PERSÖNLICHE DATEN}
\vspace{2pt}   % <-- Abstand zwischen Trennlinie und erster Zeile

\inforow{Geburtsdatum:}{21.08.2001}
\inforow{Geburtsort:}{Chennai, Indien}
\inforow{Staatsangehörigkeit:}{Indisch}
\inforow{Adresse:}{Mozartstr.\ 22, 91052 Erlangen}
\inforow{Telefon:}{+49 1522 1554845}
\inforow{E-Mail:}{\href{mailto:pravin.mani@fau.de}{pravin.mani@fau.de}}
\inforow{Sprachen:}{Deutsch (A2), Englisch (C1), Tamilisch (Muttersprache), Spanisch (A1)}

%==================ABSTANDSHALTER (zwischen Abschnitten)==================
\vspace{7pt}   % <-- Abstand zwischen Abschnitten; frei anpassbar
%=========================================================================

%==============================================================================
% AUSBILDUNG
%==============================================================================
\section{AUSBILDUNG}

\cventry{Master of Science - Physik}
        {Friedrich-Alexander-Universität Erlangen-Nürnberg}
        {Apr. 2024 - heute, \hfill Erlangen, Deutschland}

\vspace{0pt}

\cventry{Master of Science - Kernphysik}
        {Universität Madras}
        {Jul. 2021 - Aug. 2022, \hfill Chennai, Indien}

\vspace{0pt}

\cventry{Bachelor of Science - Physik}
        {Universität Madras - Gesamtnote: 8{,}96}
        {Jun. 2018 -- Apr. 2021, \hfill Chennai, Indien}

%==================ABSTANDSHALTER========================================
\vspace{7pt}
%=========================================================================

%==============================================================================
% PHYSIKPROJEKTE
%==============================================================================
\section{PHYSIKPROJEKTE}

%--- Projekt 1 ---------------------------------------------------------------
\project{Ising-Modell mit Cython}{Eigenständiges Projekt}
\begin{cvlist}
  \item Implementierung eines 2D-Ising-Monte-Carlo-Verfahrens mit C-typisierten
        Speicheransichten und Schwarz-Weiß-Metropolis-Updates in Cython.
  \item Periodische Randbedingungen und vektorfreie innere Schleifen für eine
        erhebliche Beschleunigung gegenüber reinem Python.
\end{cvlist}

\vspace{0pt}

%--- Projekt 2 ---------------------------------------------------------------
\project{In Bearbeitung -- Ising-Phasendiagramm mit verlässlicher Unsicherheit mittels JAX, FLAX}{Eigenständiges Projekt}
\begin{cvlist}
  \item Entwicklung eines FLAX-basierten CNN, das konforme Vorhersagen anwendet,
        um kalibrierte Phasenkarten mit Konfidenzdarstellungen zu erzeugen.
\end{cvlist}

\vspace{0pt}

%--- Projekt 3 ---------------------------------------------------------------
\projectdate{Vergleichende Analyse adiabatischer Fehlerschranken in Quantensystemen anhand eines modifizierten Landau-Zener-Modells}
            {Benotet: 1{,}5}
            {Dez. 2025 -- Feb. 2026}
\vspace{0pt}
\noindent\textit{FAU Erlangen-Nürnberg, bei Prof.\ Michael J.\ Hartmann, Leonhard Richter und Marco Wiedmann}
\begin{cvlist}
  \item Konstruktion eines Dreiniveau-Quantensystems und Beobachtung der simulierten
        Fehler bei kontinuierlicher und adiabatischer Entwicklung entlang der Eigenniveaus.
  \item Verwendung einer quantitativen Fehler-Oberschranke zur Verifikation der simulierten
        Fehler und Überprüfung der Schranke selbst nach Jansen, Ruskai und Seiler.
\end{cvlist}

\vspace{0pt}

%--- Projekt 4 ---------------------------------------------------------------
\project{In Bearbeitung -- Erweiterung auf diskrete Zeitentwicklung durch unitäre Operatoren und Charakterisierung des adiabatischen Fehlerverhaltens}{Eigenständiges Projekt}
\begin{cvlist}
  \item Konstruktion der unitären Operatorentwicklung mit QuTiP.
  \item Laufende Überprüfung der Fehlerschranken.
\end{cvlist}

%==================ABSTANDSHALTER========================================
\vspace{7pt}
%=========================================================================

%==============================================================================
% KENNTNISSE
%==============================================================================
\newcommand{\skillgap}{-4pt}   % <-- Abstand zwischen Kenntniszeilen; weniger negativ = mehr Abstand
\section{KENNTNISSE}
\vspace{0pt}   % <-- Abstand zwischen Trennlinie und erster Kenntniszeile
\skillrow{Programmiersprachen}{Python, C++ (Grundkenntnisse), Lisp}\vspace{\skillgap}
\skillrow{Machine-Learning-Bibliotheken}{TensorFlow, JAX, FLAX, OPTAX}\vspace{\skillgap}
\skillrow{Quantensystem-Bibliotheken}{QuTiP, QiSkit (Grundkenntnisse)}\vspace{\skillgap}
\skillrow{Entwicklerwerkzeuge}{PyCharm, Emacs, Git, Jupyter, \LaTeX}

\end{document}